 % \documentclass[10pt,respuestas,a4]{examenEPN}
\documentclass[10pt,a4]{examenEPN0}
\examen{Prueba 02}
\usepackage{comandosEPN}

\usepackage{empheq}

\newcommand{\itemcp}{\stepcounter{enumi}\item[ \textcolor{red}{\theenumi.$^*$}] }

\newcommand{\iitemcp}{\stepcounter{enumii}\item[ \textcolor{red}{\theenumii)$^*$}] }


\materia{Álgebra Lineal}
%\autor{Departamento de Formación Básica}
\autor{GR \hspace{1cm}.}
%\fecha{Semestre 2026-A}
\fecha{Nombre: \hspace{9cm} Cédula:}
\logouno[1.6cm]{Logo_EPN.eps}
\logodos[2.6cm]{Logo_DFB.eps}

\newcommand{\la}{\lambda}

%Datos
\newcommand{\datosAlgebra}[2]{
\EnExamen
{
\vspace{5mm}
    \noindent
    Nombres: \hrulefill\
    Carrera: \rule{4.5cm}{0.4pt}\
    \ifthenelse{\equal{#1}{}}%
    {Paralelo: \underline{\parbox{1.8cm}{\parbox{1cm}{}}}\\[3mm]}
    {Paralelo: \underline{\parbox{1.8cm}{\centering #1}}\\[3mm]}
    C\'edula: \rule{5cm}{0.4pt}\
    Firma: \hrulefill\
    Fecha: \underline{\parbox{2.1cm}{\centering #2}}
}}


% \newcommand{\indicacionesPrueba}{%
% \EnExamen{
% \begin{indicaciones}
%     \item La prueba tiene una duraci\'on  de \textbf{120 minutos.}
%     \item No se permite la utilizaci\'on de dispositivos electr\'onicos (celular, tablets, calculadora, etc.)
%     \item Leer detenidamente los enunciados, \textbf{justificar la respuesta para acceder al puntaje total}.
%     \item La prueba es completamente individual. Los estudiantes deberán dejar sus teléfonos celulares guardados y no podrán tenerlos consigo (por ejemplo, en los bolsillos de su ropa) durante toda la evaluación. En caso de detectarse algún acto de deshonestidad académica o el incumplimiento de esta norma, no podrán seguir rindiendo la prueba y su calificación será \textbf{cero}.
% \end{indicaciones}
% \pagestyle{empty}
% }}

\newcommand{\indicacionesPrueba}{%
\EnExamen{
\begin{indicaciones}
    \item La prueba tiene una duraci\'on  de \textbf{60 minutos.}
    \item Leer detenidamente los enunciados, \textbf{justificar la respuesta} para acceder al puntaje total.
    \item No se permite la utilizaci\'on de dispositivos electr\'onicos (celular, tablets, calculadora, etc.)
    \item La prueba es completamente individual
    \item No se permite portar teléfonos celulares, relojes inteligentes ni otros dispositivos electrónicos durante el examen, aun cuando permanezcan en los bolsillos. Estos dispositivos deberán guardarse en la mochila antes del inicio de la evaluación y esta deberá mantenerse cerrada durante todo el examen.
    \item He leído las indicaciones y acepto que el incumplimiento de estas se penaliza con \textbf{cero}. Firma:...................................

    
\end{indicaciones}
\pagestyle{empty}
}}

\newcommand{\indicacionesExamen}{
\EnExamen{\begin{indicaciones}
    \item La evaluación tiene una duraci\'on  de \textbf{120 minutos.}
    \item No se permite la utilizaci\'on de dispositivos electr\'onicos (celular, tablets, calculadora, etc.)
    %\item Leer detenidamente los enunciados, justificar la respuesta para acceder al puntaje total.
    \item El examen es completamente individual. En caso de detectarse alg\'un acto de deshonestidad acad\'emica no podr\'a seguir rindiendo la prueba y su calificaci\'on ser\'a \textbf{cero}.
\end{indicaciones}
}\pagestyle{empty}
}

%% Imprimir las una hoja A4, dos veces en A5
\newcommand{\imprimirDoble}
{%
\usepackage{pgfmorepages}
\pgfmorepagesloadextralayouts
\pgfpagesuselayout{repeated 2-up}[a4paper,landscape]
\thispagestyle{empty}
}


%Este sí se puede comentar
\newenvironment{solucion}{\color{blue}\proof[Solución]}{\endproof\color{black}}




% \renewcommand{\@autor}{
% ---\\
% \phantom{.} 
% ...}
% \makeatother

\imprimirDoble

\begin{document}
\encabezado
\datosAlgebra{GR5M}{04/11/2026}
%Poner el grupo y la fecha, o dejar las llaves vacías.
\indicacionesPrueba



\begin{preguntas}

%%%%%%%%%%%%%%%%%%%%%%%%%%%%%%%%%
%%%%%%%%%%%%%%%%%%%%%%%%%%%%%%%%%

% \item Se desea preparar un fertilizante que contenga tres tipos de abono. El abono A: no contiene Nitrógeno,
% contiene 2 kg de Fósforo y 1 kg de Potasio por saco. El abono B: contiene 1 kg de Nitrógeno, 1 kg de Fósforo y 1 kg de Potasio por saco. El abono C: contiene 2 kg de Nitrógeno, 3 kg de Fósforo y 1 kg de Potasio por saco. La mezcla del fertilizante requerido debe contener 7 kg de Nitrógeno, 14 kg de Fósforo y 6 kg de Potasio.

\item Se desea preparar un fertilizante que contenga tres tipos de abono. 

\begin{itemize}
\item El abono A: no contiene Nitrógeno,
contiene 2 kg de Fósforo y 1 kg de Potasio por saco. 
\item El abono B: contiene 1 kg de Nitrógeno, 1 kg de Fósforo y 1 kg de Potasio por saco. 
\item El abono C: contiene 2 kg de Nitrógeno, 3 kg de Fósforo y 1 kg de Potasio por saco. 
\end{itemize}
La mezcla del fertilizante requerido debe contener 7 kg de Nitrógeno, 14 kg de Fósforo y 6 kg de Potasio.

\begin{enumerate}
\item Defina las tres variables, plantee el sistema de ecuaciones lineales que modela la situación y escriba su representación matricial $Ax=b$, identificando $A$, $x$ y $b$. \puntaje{3}
\item Resuelva el sistema mediante el método de la matriz inversa y determine los valores de las tres incógnitas en el contexto del problema. \puntaje{3}
\end{enumerate}

%%%%%%%%%%%%%%%%%%%%%%%%%%%%%%%%%%%%%%%
\item De un ejemplo de una matriz diagonal $A\in \mathbb{R}^{3\times 3}$ y $b\in \mathbb{R}^{3\times 1}$ que satisfacen la ecuación del sistema matricial $Ax= b$, si se conoce que \puntaje{6}
        \[
            \mathrm{det}(A) = 3,
            \quad 
            \mathrm{det}(A^{(1)}) = -9, 
            \quad
            \mathrm{det}(A^{(2)}) = 12, 
            \quad
            \mathrm{det}(A^{(3)}) = 6, 
            \quad
        \]
        
%%%%%%%%%%%%%%%%%%%%%%%%%%%%%%%%%%%%%%%
% \item (Ejercicio del deber)    Considere el sistema de ecuaciones lineales: \puntaje{6}
% $$
% \begin{cases} x + 2y + z = 4 \\ 2x - y + z = 2 \\ 3x + y - 2z = 2 \end{cases}
% $$
% Determine la solución $(x, y, z)$ empleando la Regla de Cramer.

\item (Ejercicio del deber)  Determine la solución $(x, y, z)$ del siguiente sistema empleando la Regla de Cramer.  \puntaje{6}
$$
\begin{cases} x + 2y + z = 4 \\ 2x - y + z = 2 \\ 3x + y - 2z = 2 \end{cases}
$$


%%%%%%%%%%%%%%%%%%%%%%%%%%%%%%%%%%%%%%%%%
\item  Sean $A$ y $B$ dos matrices de orden $n \times n$ tales que son semejantes (es decir, existe una matriz invertible $P$ tal que $A = P B P^{-1}$). Si $\det(A) = 5$, ¿cuál es el valor de $\det(B)$? 
Justifique su respuesta \puntaje{2}
    \begin{enumerate}[label=(\alph*), itemsep=2pt]
        \item $\dfrac{1}{5}$
        \item $5$
        \item $-5$
        \item Depende del valor de $\det(P)$.
    \end{enumerate}   
    
\end{preguntas}
\final
\end{document}